\label{sec:implementation}
\subsection{Generation of Feynman diagram expressions}

We generate the Feynman diagrams including symmetry factors and signs
(from closed fermion loops) with the help of the program
\texttt{qgraf}\,\cite{Nogueira:1991ex,Nogueira:2006pq}. In the notation
of this program, the propagators for the gluon, the ghost, and a single
quark flavor can be defined as
\begin{verbatim}
[g,g,+], [c,C,-], [fq,fQ,-].
\end{verbatim}
The sign in the third entry of each square bracket
denotes whether the particle is a boson or a fermion.

The flow lines for the gluon and the quark are implemented by
introducing separate fields \texttt{b}, \texttt{fr}, and \texttt{fs} for
$L$, $\lambda$, and $\bar\lambda$, respectively. They are implemented as
\begin{verbatim}
[b,B,+], [fr,fR,-], [fs,fS,-],
\end{verbatim}
where the latter two represent the $\langle\chi\bar\lambda\rangle$ and
$\langle\lambda\bar\chi\rangle$ bilinears (see \eqs{eq:chibarlambda} and
\noeqn{eq:lambdabarchi}). With these fields, we can then define the
regular as well as the flow vertices. For example, for the trilinear
flow vertex of the pure gauge theory defined in \eqn{eq:3gvert}, we
define
\begin{verbatim}
[B,b,b], [B,b,g], [B,g,g],
\end{verbatim}
which corresponds to the three combinations of
\fig{fig:3gcomb}.

Already in regular \qcd\ it is convenient to separate the color
structure from the rest of the calculation. This becomes non-trivial in
cases which involve the four-gluon vertex. It is thus convenient to
introduce an auxiliary ``particle'' $\Sigma_{\mu\nu}^{a}$ whose
``propagator'' is given by\,\cite{Chetyrkin-priv}
\begin{align}
  &
  \begin{gathered}
    \includegraphics{images/sigma_propagator.pdf}
  \end{gathered}
  & = &
  \begin{split}
    \delta^{ab}\delta_{\mu\rho}\delta_{\nu\sigma} \,.
  \end{split}
  & &
  %% \label{eq:}
\end{align}
The four-gluon vertex can then be replaced by a trilinear $gg\Sigma$
vertex whose Feynman rule reads
\begin{align}
  &
  \begin{gathered}
    \includegraphics{images/sigma_vertex.pdf}
  \end{gathered}
  & = &
  \begin{split}
    \frac{\mathrm{i}g}{\sqrt{2}} f^{abc}\left(\delta_{\mu\rho}\delta_{\nu\sigma}-\delta_{\nu\rho}\delta_{\mu\sigma}\right) \,.
  \end{split}
  & &
  %% \label{eq:}
\end{align}
This allows to factorize the color factor off all Feynman diagrams,
albeit at the cost of increasing their number. We proceed
correspondingly for the quartic gluon flow vertices by introducing an
additional $\Sigma$ particle which is directed in flow time. It forms
trilinear vertices with gluons and gluon flow lines, as shown
explicitly in Appendix\,\ref{sec:frules}. Again, they contain
exactly one outgoing ($\Sigma$ or gluon) flow line.

By default, \texttt{qgraf} also generates diagrams with closed flow-line
loops.  We use a \texttt{Perl} script to parse the \texttt{qgraf} output
in order to eliminate such diagrams before the actual calculation. They
vanish trivially algebraically though, as discussed earlier.  In
the same way we multiply diagrams by a factor $\NF$ for each closed
fermion loop.  We then process the diagrams with the help of
\texttt{q2e/exp}\,\cite{Harlander:1997zb, Seidensticker:1999bb} in order
to insert the Feynman rules and convert the diagrams to \texttt{FORM}
code.

Within \texttt{FORM}\,\cite{Vermaseren:2000nd, Kuipers:2012rf}, we
contract the Lorentz indices, take the fermion traces, and simplify the
resulting expressions to a standard form of scalar integrals, as defined below.
  The color factor is evaluated separately with the
help of the \texttt{color} package\,\cite{vanRitbergen:1998pn}.


For example at the three-loop level (as needed for the results in this paper)
the flow-time integrals can be described in the following form
\begin{equation}
\begin{split}
&I(t,\mathbf{n},\mathbf{a},\mathbf{c},D) =
  \left(\prod_{r=1}^N\int_0^{t_{r}^\text{up}}\dd
  t_r\,t_r^{c_r}\right)\int_{p_1,p_2,p_3} \frac{\exp[\sum_{k,i,j}
      a_{kij} t_k p_i \cdot p_j] }{p_1^{2n_1}p_2^{2n_2}p_3^{2n_3}
    p_4^{2n_4}p_5^{2n_5}p_6^{2n_6}}\,,
\label{eq:intrep}
\end{split}
\end{equation}
where
\begin{equation}
\begin{split}
\mathbf{n} &= \{n_1,\ldots,n_6\}\,,\qquad\mathbf{c} = \{c_1,\ldots,c_N\}\,,\\
\mathbf{a} &= \{a_{kij}:k=0,\ldots, N\,;\,i=1,2,3\,;j=1,2,3\}\,,\\
\end{split}
\end{equation}
are sets of integers, $N\leq 4$, $t_0\equiv t$, and the upper limits for
the flow-time integrations are linear combinations of the other
flow-time variables, $t_r^\text{up}=t_r^{\text{up}}(t_0,\ldots,
t_{r-1})$. Initially, all the $c_i$ are zero. It is helpful to introduce
these parameters for the subsequent discussion though.  The momenta
$p_4$, $p_5$, $p_6$ are linear combinations of the integration momenta
$p_1$, $p_2$, $p_3$.  For the three-loop quantities considered
in this paper, the flow time $t$ is the only dimensionful external scale.
A generalization of this notation to a different number of loops and to additional external
scales is straightforward.

\subsection{\abbrev{IBP} reduction of flow-time loop integrals}\label{sec:reduction}

In the first step, we reduce all occurring integrals to so-called master
integrals by employing integration-by-parts (\abbrev{IBP}) identities\,\cite{Tkachov:1981wb,Chetyrkin:1981qh}.
They are based on the
observation that integrals over total derivatives vanish in dimensional
regularization. Applying the operator
\begin{equation}
  \begin{split}
    \mathcal{D}_{ij} \equiv \frac{\partial}{\partial p_i}\cdot p_j =
    \delta_{ij}\,D+ p_j\cdot \frac{\partial}{\partial p_i}
    \label{eq:dij}
  \end{split}
\end{equation}
to the integrand of $I(t,\mathbf{n},\mathbf{a},\mathbf{c},D)$ in
\eqn{eq:intrep} thus results in a vanishing integral for all
$i,j\in\{1,2,3\}$. On the other hand, explicitly acting with the
derivative on the integrand leads to a sum of integrals where some of the
indices $n_i$ and $c_r$ are shifted by $\pm 1$.\footnote{To be
  precise: At most one $n_i$ is shifted by $-1$ and at most one of the
  $n_i$ and at most one of the $c_r$ are shifted by $+1$.} Repeated
application of the $\mathcal{D}_{ij}$ thus leads to linear relations
among integrals with different $\mathbf{n}$ and $\mathbf{c}$.
However, the number of flow-time integrations and the exponential
function in the integrand still remain unaltered by this procedure.

We may apply an analogous strategy for the flow-time parameters though.
Inserting a derivative with respect to one of the flow-time integration
variables, one arrives at the sum of two integrals with fewer flow-time
integrations and an altered exponential,
\begin{equation}
  \int_0^{t_{r}^\text{up}}\dd t_r \deriv{}{t_r} f(t_r,\ldots) = f(t_{r}^\text{up},\ldots)-f(0,\ldots)\,.
  \label{eq:ibpflow}
\end{equation}
On the other hand, explicit evaluation of the derivative at the
integrand level either reduces one of the indices $c_k$ or one of the indices $n_i$
by one. One therefore arrives at linear relations among integrals with
different $\mathbf{n}$, $\mathbf{c}$, $\mathbf{a}$, and different
number of flow-time integrations $N$.

Applying the above operations to ``seed integrals'', i.e., integrals
with fixed numerical values for the $\mathbf{n}$, $\mathbf{c}$, and
$\mathbf{a}$, allows one to build a system of linear relations among the
$I(t,\mathbf{n},\mathbf{a},\mathbf{c},D)$. Defining an ordering
(complexity) criterion allows one to solve the system using a Gaussian elimination
type reduction for a minimal set of integrals (master integrals) \cite{Laporta:2001dd}.
Through this procedure all integrals of the type defined in \eqn{eq:intrep}
are reduced to the minimal set of master integrals, which are simplest
by the ordering criterion introduced. The system of linear relations
is process dependent and for the observables in this study we construct it
in \texttt{Mathematica}\,\cite{Mathematica11.3}. We use the specialized software
\texttt{Kira}\,\cite{Maierhoefer:2017hyi,Maierhofer:2018gpa} for the reduction of
such linear equation systems to solve it.

While \texttt{Kira} alone is sufficient
for the reduction of all our integrals at the two-loop level (see
e.g.\ \citere{Harlander:2018zpi}), we find that the algebraic solution
of the system at the three-loop level would require more than
750\,GiB of \abbrev{RAM} and thus exceeds our available computing
resources.\footnote{This calculation was performed before the release of
  \texttt{Kira 1.2}\,\cite{Maierhofer:2018gpa}. It is well possible that
  these statements could change with the newly implemented features.}
Using finite field and
reconstruction techniques, which
over the last decade have gained an increased use for higher-order
perturbative calculations (see,
e.g.\ \citeres{Kauers:2008zz,Kant:2013vta,Peraro:2016wsq}), one can improve on the required computational resources as
follows.

Since $t$
is the only dimensionful scale and can thus be factored out, the
coefficients of the integrals in the linear system are simply
polynomials in $D$. The individual steps for the reduction of the linear system only involve
elementary arithmetic operations, which means that the coefficients of the master
integrals are rational functions in $D$. One can then solve the linear system
numerically over a finite field using a sufficiently large number of different integer
values for $D$ (``probes'') and reconstruct the rational functions in $D$ exactly. With this approach
the size of the coefficients remain simple numbers and the requirements on computational
resources can be improved.

While \texttt{Kira} already uses \texttt{pyRed} as a first step to remove
linearly dependent equations over a finite field, in our approach \texttt{pyRed} is used
to reduce the system itself multiple times over a finite field.
The resulting numbers are processed with the
library \texttt{FireFly}\,\cite{Klappert:2019emp}, which provides an
efficient implementation of interpolation\,\cite{Zippel:1990,Cuyt:2011}
and rational-reconstruction
algorithms\,\cite{Wang:1981,Monagan:2004}.\footnote{We remark that
  \texttt{FireFly} also works for multi-variate rational functions.}  In
our case, we require 201 probes, chosen from three different finite
fields, defined as prime fields $\mathbb{Z}_p$ with $p$ a 63-bit prime
number, to reconstruct all coefficients, plus one additional probe in a
fourth prime field to verify the reconstruction. For each probe, the
solution of the system with \texttt{Kira} now just requires about 70\,GiB of
\abbrev{RAM} and takes about three \abbrev{CPU} hours.  In total, the reduction took
less than three days, using ten threads on two Intel Xeon Gold 6138
processors.

For the observables considered further below in this paper at the three-loop level we start
with a total of 3195 integrals and can reduce them to a minimal set of 188 master integrals.\footnote{When our
	observables are expressed in terms of the master integrals, the dependence of six of them drops out.}
Their numerical evaluation is described in the next section.

\subsection{Numerical computation of flow-time loop integrals}

As a first step of the numerical evaluation of the gradient-flow integrals given by \eqn{eq:intrep},
we express the propagators through Schwinger-parameter integrals and map them from $x\in(0,\infty)$ to
$y\in(0,1)$ using the simple transformation $x=y/(1-y)$.
Momentum integrations are performed as $D$-dimensional Gaussian integrals after a diagonalization.
Similarly, all
flow-time integrations are mapped to the unit interval with simple
linear transformations. The overall result is an integral over a unit hypercube.

The integrand typically involves a number of (overlapping)
singularities, which we factorize using
\texttt{FIESTA}\,\cite{Smirnov:2013eza}, an implementation of the sector
decomposition algorithm\,\cite{Binoth:2003ak}.  As opposed to regular
Feynman integrals where sector decomposition is typically employed in
combination with Feynman parameterization, it appears that we cannot
restrict the singularities to the lower integration bound
only.\footnote{To identify cases where singularities appear at the upper
  bound, we determine the degree of divergence for each subset of
  integration variables at the level of the Schwinger parameterization
  \cite{Panzer:2014gra}.} If an integral involves singularities both at the
lower and the upper bound, we split all integration intervals in the
middle, and map the singularities to the lower bound by an appropriate
change of variable.

The sector decomposed integrals obtained from \texttt{FIESTA} are then
integrated with our implementation of fully symmetric integration rules
of order 13\,\cite{GenzMalik,Harlander:2016vzb}, which can handle
integrable logarithmic-like singularities at the integration boundaries
very well. We perform all arithmetic with 256-bit precision using the
\texttt{MPFR} library\,\cite{Fousse:2007}, and used a local adaptive
bisection in the direction of the largest fourth difference. A high
precision arithmetic turns out to be necessary to achieve a relative
numerical accuracy of $10^{-10}$ or better. The integration uncertainty
is estimated by the difference between the integration results of rules
of order 13 and 11.

As a check of our results further below, we apply the numerical integration method
to our unreduced set of $3195$ integrals as well as to the set of $188$ master integrals.
Finding full agreement for the results obtained with both sets within numerical uncertainties
serves as a check of the numerical integration as well as for the reduction procedure.

\paragraph{Analytical computation.}
For some master integrals we can find analytical results either
through elementary methods with the help of
\texttt{Mathematica}\,\cite{Mathematica11.3} or by using
\texttt{HyperInt} \cite{Panzer:2014caa}.
The integration with \texttt{Mathematica} sometimes yields hypergeometric functions which can be expanded with the package \texttt{HypExp}\,\cite{Huber:2005yg,Huber:2007dx}.
For example as shown further below, we
find analytical results for all contributions which contain at least one factor
of $\TR$ in the color structure. A fully
analytical calculation of integrals with other color factors requires a more
detailed investigation. However, for all practical purposes, the
numerical results provided in this paper are (more than) sufficient.
